\documentclass[11pt,a4paper]{article}
\usepackage[T1]{fontenc}
\usepackage[utf8]{inputenc}
\usepackage{amsmath,amssymb,amsthm}
\usepackage[margin=1in]{geometry}
\usepackage[colorlinks=true,linkcolor=blue,urlcolor=blue]{hyperref}
\theoremstyle{plain}
\newtheorem{theorem}{Theorem}
\newtheorem{lemma}[theorem]{Lemma}
\newtheorem{proposition}[theorem]{Proposition}
\newtheorem{corollary}[theorem]{Corollary}
\theoremstyle{definition}
\newtheorem{remark}[theorem]{Remark}

\title{Recovering the unwritten recurrences of the table OEIS A251801}
\author{Adrian Perez Fontelles\\ \small Independent researcher}
\date{2 September 2026}
\begin{document}
\maketitle

\begin{abstract}
OEIS A251801 is a two-dimensional table $T(n,k)$ whose empirical recurrences are, for several of
its lines, given only as an ORDER --- ``k=2: [order 20]'' and the like --- with the
coefficients in a linked file rather than in the entry. Those conjectures can still be settled,
and the recurrences recovered. A line of the table is a fixed-width or fixed-height array
count, hence a walk count on finitely many vertices, hence satisfies some monic recurrence of
order at most that number; Berlekamp--Massey on twice as many exact terms returns the MINIMAL
one. Where its order equals the order the entry states --- and where the same holds after the
first few terms are dropped --- any recurrence of that order the line satisfies has a
characteristic polynomial that is a multiple of the minimal one and of equal degree, hence is
that polynomial. This note settles column 2, row 2, row 3.
\end{abstract}

\noindent\small 2020 Mathematics Subject Classification. 05A15, 11B37, 15A18.\normalsize

\section{The table and the unwritten conjectures}

OEIS A251801 is ``T(n,k)=Number of (n+1)X(k+1) 0..3 arrays with no 2X2 subblock having x11-x00 less than x10-x01.''. It is read by antidiagonals and begins
\[
150, 1080, 1740, 6627, 28932, 19269, 36552, 348679,\ \dots
\]
The entry carries, among its empirical claims:
\begin{quote}\small
k=2: [order 20]\\
n=2: [order 22]\\
n=3: [order 40]
\end{quote}
As of the ``Last modified'' line on the live entry (Jun 02 2025 10:32:19, revision 6) these are still
recorded as empirical, and the recurrences themselves are not in the entry text.

\section{Why a line of the table is a walk count}

Fix the index that the line fixes. The entry defines $T(n,k)$ as a number of arrays of a shape
determined by $n$ and $k$, subject to a condition local to a bounded window of consecutive
lines. Substituting the fixed index into the entry's own wording turns the definition into an
ordinary one-parameter array count, and for such a count the admissible windows are the
vertices of a finite digraph and an array is a walk. So the line is a walk count on $S$
vertices, and by the Cayley--Hamilton theorem it satisfies a monic integer recurrence of order
at most $S$.

\section{Recovering the recurrences}

\begin{lemma}\label{lem:bm}
A sequence known to satisfy some linear recurrence of order at most $S$ is determined, with its
minimal recurrence, by its first $2S$ terms; Berlekamp--Massey applied to them returns that
minimal recurrence exactly.
\end{lemma}

\subsection*{Column $k=2$}
Substituting the fixed index into the entry's wording gives an ordinary one-parameter array
count, a walk on $S=64$ vertices. Berlekamp--Massey on $2S=128$ exact terms
returns a minimal recurrence of order $20$ --- the order the entry states --- with
integer coefficients
\begin{center}\small\begin{tabular}{cccccccc}
$c_{1}$ & $30$ & $c_{6}$ & $-170647$ & $c_{11}$ & $2262184$ & $c_{16}$ & $-12352$ \\
$c_{2}$ & $-40$ & $c_{7}$ & $-169762$ & $c_{12}$ & $3737187$ & $c_{17}$ & $-613952$ \\
$c_{3}$ & $-2202$ & $c_{8}$ & $1201353$ & $c_{13}$ & $-3793360$ & $c_{18}$ & $105792$ \\
$c_{4}$ & $7167$ & $c_{9}$ & $-237196$ & $c_{14}$ & $-1475516$ & $c_{19}$ & $54144$ \\
$c_{5}$ & $41092$ & $c_{10}$ & $-3305290$ & $c_{15}$ & $2385944$ & $c_{20}$ & $-13824$ \\
\end{tabular}\end{center}
so that $a(m)=\sum_{i=1}^{20}c_i\,a(m-i)$ for every $m$. Removing the first
$20+4$ terms and repeating gives order $20$ again, which is what the
identification below needs.

\subsection*{Row $n=2$}
Substituting the fixed index into the entry's wording gives an ordinary one-parameter array
count, a walk on $S=64$ vertices. Berlekamp--Massey on $2S=128$ exact terms
returns a minimal recurrence of order $22$ --- the order the entry states --- with
integer coefficients
\begin{center}\small\begin{tabular}{cccccccc}
$c_{1}$ & $40$ & $c_{7}$ & $8833486$ & $c_{13}$ & $492639700$ & $c_{19}$ & $20037984$ \\
$c_{2}$ & $-755$ & $c_{8}$ & $-27833039$ & $c_{14}$ & $-498877265$ & $c_{20}$ & $-4463424$ \\
$c_{3}$ & $8946$ & $c_{9}$ & $71958408$ & $c_{15}$ & $414825042$ & $c_{21}$ & $625536$ \\
$c_{4}$ & $-74677$ & $c_{10}$ & $-153781873$ & $c_{16}$ & $-280100140$ & $c_{22}$ & $-41472$ \\
$c_{5}$ & $467156$ & $c_{11}$ & $272810702$ & $c_{17}$ & $151065320$ &  &  \\
$c_{6}$ & $-2274363$ & $c_{12}$ & $-402324879$ & $c_{18}$ & $-63500432$ &  &  \\
\end{tabular}\end{center}
so that $a(m)=\sum_{i=1}^{22}c_i\,a(m-i)$ for every $m$. Removing the first
$22+4$ terms and repeating gives order $22$ again, which is what the
identification below needs.

\subsection*{Row $n=3$}
Substituting the fixed index into the entry's wording gives an ordinary one-parameter array
count, a walk on $S=256$ vertices. Berlekamp--Massey on $2S=512$ exact terms
returns a minimal recurrence of order $40$ --- the order the entry states --- with
integer coefficients
\[
(c_1,\dots,c_{40})=(70,\ -2375,\ 52042,\ -827990,\ 10195132,\ -101124322,\ 830421108,\ -5758451859,\ 34223487178,\ -176332070653,\ 794779755110,\ -3156457476820,\ 11109873120920,\ -34818233254780,\ 97527763758712,\ -244892909297175,\ 552543768062970,\ -1122181309818449,\ 2053995521193590,\ -3390759101221990,\ 5049772850347484,\ -6783271969135666,\ 8213149434578420,\ -8953247189169461,\ 8772430272272310,\ -7708025280163467,\ 6055906423103258,\ -4238615566223888,\ 2630732733430288,\ -1439628171650384,\ 689679083321824,\ -286667447021696,\ 102211164419072,\ -30802911196416,\ 7693171914240,\ -1549479297024,\ 241805316096,\ -27432235008,\ 2012553216,\ -71663616).
\]
so that $a(m)=\sum_{i=1}^{40}c_i\,a(m-i)$ for every $m$. Removing the first
$40+4$ terms and repeating gives order $40$ again, which is what the
identification below needs.

\begin{theorem}
For each line above, the displayed recurrence holds for every index, and it is the recurrence
the entry records.
\end{theorem}

\begin{proof}
It holds by Lemma~\ref{lem:bm}. For the identification, the entry asserts that the line
satisfies SOME recurrence of the stated order, possibly only from some index on. Let $q$ be its
characteristic polynomial and $q_{\min}$ the minimal polynomial of the tail. Then
$q_{\min}\mid q$, and the second Berlekamp--Massey run --- on the line with its first few
terms removed --- shows $\deg q_{\min}$ equals the stated order, which is $\deg q$. Both are
monic, so $q=q_{\min}$.
\end{proof}

\section{Verification}

Three checks.

First, each line's model was compared against the entry's own published values for that line,
taken from the antidiagonal DATA read in either orientation and from the ``Table starts''
block, and agrees at every available term. This is the only point at which reading the entry's
English enters, and a misreading gives different counts.

Second, each recovered recurrence was evaluated directly on the model's values, with no
Berlekamp--Massey involved, and holds at every index tested.

Third, each order was computed twice by independent routes --- modulo a $61$-bit prime, where
every intermediate is a machine word, and again in exact rational arithmetic --- and once more
on the line with its first few terms dropped, which is what the identification requires. All
agree.

\begin{thebibliography}{9}
\bibitem{oeis} The OEIS Foundation, \emph{The On-Line Encyclopedia of Integer
Sequences}, \url{https://oeis.org}, 2026. Sequence A251801.
\bibitem{massey} J.~L.~Massey, Shift-register synthesis and BCH decoding, \emph{IEEE
Trans. Inform. Theory} \textbf{15} (1969), 122--127.
\bibitem{stanley} R.~P.~Stanley, \emph{Enumerative Combinatorics}, Volume 1, 2nd ed.,
Cambridge University Press, 2011. (Transfer-matrix method, Section 4.7.)
\bibitem{horn} R.~A.~Horn and C.~R.~Johnson, \emph{Matrix Analysis}, 2nd ed., Cambridge
University Press, 2013.
\end{thebibliography}

\end{document}
